\documentclass[11pt]{article}
\usepackage[a4paper,margin=27mm]{geometry}
\usepackage[T1]{fontenc}
\usepackage{lmodern,microtype,amsmath,amssymb,amsthm,enumitem}
\usepackage[hidelinks]{hyperref}
\newtheorem{theorem}{Theorem}
\newtheorem{lemma}[theorem]{Lemma}
\newtheorem{corollary}[theorem]{Corollary}
\theoremstyle{remark}\newtheorem{remark}[theorem]{Remark}
\newcommand{\Dscr}{\mathcal D}
\newcommand{\cred}{\mathrm{credit}}
\newcommand{\debt}{\mathrm{debt}}
\newcommand{\J}{\mathcal J}
\newcommand{\K}{\mathcal K}
\newcommand{\U}{\mathcal U}
\title{Fineness for Any Finite Number of Worlds}
\author{}\date{11 October 2026}
\begin{document}\maketitle

\noindent
This note answers the higher-dimensional propagation questions in
\emph{Fineness Spreads Along Edges} (FSE). Under its complete-ranking
hypotheses, one fine world makes every world fine for every finite
$n\ge3$. A fine edge, or a world with a one-sided finitely additive
extension, also forces the vacuous model to be dominated in every such
dimension. Its conclusions that no non-null region is credit-null or
debt-infinitesimal also hold for every finite $n\ge3$.
The proofs use three coordinates at a time, without restricting the
rankings or the selection to a three-world problem. The remaining
coordinates are left unrestricted. These arguments are author-checked,
not independently refereed.

\section{Setting and the inputs from FSE}

There are finitely many worlds, $n\ge3$. We use FSE's measurable models
modulo Gaussian null sets, guidance ranks $0,1,2,3$ (good, silent, bad,
contradictory), and regular models (correct on unanimous problems and
never contradictory). Each world's ranking is a complete preorder,
with smaller better, satisfying strict guidance improvement (G) and
independence under replacement of common advice (IB).
The selection axiom (S) supplies coherent undominated $C_p$, competitive
on $\operatorname{supp}p$, with every section
$\{p:C_p\preceq_iX\}$ closed. Write $V=\Dscr_\Delta$.

Let $H'$ be the mixed binary problems. For distinct $i,j$ put
\[
 D_{ij}=\{g:g_ig_j<0\},\qquad t^i_j=|g_i|/|g_j|\quad\hbox{on }D_{ij}.
\]
Problems may be represented in the half-space $g_i>0$ when working at
world $i$; changing representative by $g\mapsto-g$ preserves all advice
and null sets. Credit and debt mean good and bad advice at the world in
question. For measurable sets define
\begin{align*}
 A\preceq^i_dB&\iff V+\debt_i(A)\preceq_iV+\debt_i(B),\\
 A\preceq^i_cB&\iff V+\cred_i(B)\preceq_iV+\cred_i(A).
\end{align*}
These comparisons transfer to every common context by (IB), including
when $A,B$ overlap. Strict inclusion is strict in both orders by (G).
For non-null $A,T$, let $\J^i_A$ consist of the subsets of finite unions
of sets $\preceq^i_d A$, and define $\K^i_T$ similarly for credit.
They are ideals. World $i$ is \emph{fine} if $H'\in\J^i_A\cap\K^i_T$
for every non-null $A,T$.

A disjoint debt/credit pair $(A,T)$ is SI at $i$ if
$V+\debt_i(A)+\cred_i(T)\prec_iV$, and strictly AC if the reverse
strict comparison holds. We use the following results of FSE, each of
which is already proved there for arbitrary finite $n\ge3$:
\begin{enumerate}[label=(\arabic*),leftmargin=*,itemsep=3pt]
\item On an edge $[e_i,e_j]$, a selection has thresholds
$0\le\beta\le\alpha\le\infty$: debt at $i$ where $t^i_j<\beta$,
silence between the thresholds, and credit where $t^i_j>\alpha$.
There are edge indices $q_m\to q_e$ with $C_{q_e}=D_j$,
$C_{q_m}\ne D_j$, and $\beta_m\to\infty$ (FSE, vertex-block lemma).
\item For any non-null $A$, on every edge at $i$ a sufficiently high
tail is strictly smaller than $A$ in the debt order. For any non-null
$T$, a sufficiently low cone is no larger than $T$ in the credit order.
There are non-null debt and credit yardsticks at any prescribed bounded
interval of edge levels (FSE, yardstick lemma).
\item Along the sequence in (1), if $\alpha_m<\infty$, then for every
non-null $A'$ at bounded $t^i_j$-levels, $(A',W)$ is strictly AC for all
large $m$ and every non-null $W\subseteq\{t^i_j>\alpha_m\}\cap D_{ij}$
(FSE, small-AC-pair lemma).
\item Fineness on $D_{ij}$ transfers between its two worlds. If it holds,
then $D_{ij}\in\J^i_A\cap\K^i_T$ for all non-null $A,T\subseteq H'$.
\item A fine world on $D_{jk}$ has a two-sided local extension there:
\[
 \Lambda_j(M)=P_d(\debt_j(M)\cap D_{jk})-
 \kappa P_c(\cred_j(M)\cap D_{jk}),\qquad 0<\kappa<\infty.
\]
The charges vanish exactly on null sets. Strict inequalities of
$\Lambda_j$ give strict comparisons for regular pairs agreeing off
$D_{jk}$. For each $\epsilon>0$, $D_{jk}$ has a finite partition into
pieces of both $P_d$- and $P_c$-mass less than $\epsilon$.
\item If the partner on an edge is fine, every non-null credit yardstick
has an SI partner: for every non-null $T$ there are non-null $T_y$
inside that edge's disagreement region, $T_y\preceq^i_cT$, and $D'$ such
that $(D',T_y)$ is SI at $i$ (FSE, SI-partner lemma).
\end{enumerate}
Inputs (4)--(6) are used only where their hypotheses are explicitly
available. We do not assume fineness elsewhere.

An ultrafilter of the measure algebra contains no null set, chooses one
member of each finite partition, and is closed under intersections and
supersets. If $R\notin\J^i_A$, some ultrafilter contains $R$ and avoids
$\J^i_A$; the credit version is the same (FSE, ultrafilter lemma).
Completeness then gives $A\prec^i_dB$ for every $B$ in that ultrafilter.
If it contains $D_{ij}$, it contains all one-sided bands of a unique
level $z\in[0,\infty]$. Such a level cannot be $\infty$ for a debt-heavy
ultrafilter, or $0$ for a credit-heavy one, by (2).

\section{The part not covered by the fine world}

Fix three distinct worlds $i,j,k$. The region
\[
 R= D_{ij}\setminus D_{ik}
   =\{g:g_ig_k>0,\ g_ig_j<0\}
\]
lies in both $D_{ij}$ and $D_{jk}$. At these problems $i,k$ agree and
$j$ disagrees. Every nonempty interval of $t^i_j$ has a non-null part
in $R$. This remains true in higher dimensions, because no other
coordinate is constrained.

\begin{lemma}\label{lem:remaining}
Suppose world $j$ is fine on $D_{jk}$.
\begin{enumerate}[label=(\alph*),leftmargin=*]
\item No ultrafilter containing $R$ avoids $\J^i_A$ for a non-null $A$.
\item Suppose world $i$'s debt order is fine. If, for every non-null $T$,
there are $T_y\preceq^i_cT$ disjoint from $R$ and an SI pair $(D',T_y)$
with both parts non-null, no ultrafilter containing $R$ avoids $\K^i_T$.
\end{enumerate}
\end{lemma}
\begin{proof}
Suppress $i$ from the debt and credit orders and put $t=t^i_j$.
Take the sequence $C_m=(\alpha_m,\beta_m)$ in input (1).
On $R$, credit at $i$ is debt at $j$, and debt at $i$ is credit at $j$.
Thus world $j$'s local extension gives the following increments, for
changes supported in $R$:
\[
 \debt_i\to\text{silence}:a(B),\qquad
 \text{silence}\to\cred_i:b(B),\qquad a=\kappa P_c,\quad b=P_d.
\]
If $\U$ contains $R$, it contains subsets of any prescribed member of
arbitrarily small $a+b$ mass, by the finite partitions in input (5).
All applications of $\Lambda_j$ below change advice only inside $R$.

(a) Let $\U$ avoid $\J_A$ and contain $R$. Its level $z$ is finite.
For large $m$, $R\cap\debt_i(C_m)\in\U$. Pass to a subsequence with
either $\alpha_m=\infty$ throughout or $\alpha_m<\infty$ throughout.

In the first case fix a large $m$. By (2) a non-null sufficiently high
tail $B'\subseteq R\cap\text{silence}(C_m)$ has $B'\preceq_d A$.
Choose $B\in\U$ inside $R\cap\debt_i(C_m)$ with $a(B)<a(B')$.
Heaviness gives $B'\prec_d B$. Moving debt from $B$ to $B'$ therefore
strictly improves $C_m$ at $i$, by (IB), and changes $\Lambda_j$ by
$a(B)-a(B')<0$. This improves it at both worlds of its competitive
support, a contradiction.

In the second case choose a non-null yardstick $A'\preceq_d A$ at
bounded $t$-levels, and fix large $m$. Take $m'$ still larger with
$\alpha_{m'}>\alpha_m$, and put $W=R\cap\{t>\alpha_{m'}\}$.
Then $(A',W)$ is strictly AC by (3), and $W\subseteq\cred_i(C_m)$.
Since $A'\in\J_A$, choose $B\in\U$ inside
$R\cap\debt_i(C_m)\setminus A'$ with $a(B)<b(W)$.
Now $A'\prec_d B$, so $(B,W)$ is strictly AC by (IB).
Silencing $B\cup W$ in $C_m$ strictly improves it at $i$, and changes
$\Lambda_j$ by $a(B)-b(W)<0$, again a contradiction.

(b) Let $\U$ contain $R$ and avoid $\K_T$. Choose $T_y,D'$ as in the
hypothesis. Then $\U$ avoids $\K_{T_y}$, so $T_y\prec_c B$ for every
$B\in\U$. Choose a proper non-null part $D'_0$ of $D'$.
By debt fineness, $R\cap\{t<1\}$ has a non-null subset $A'$ with
$A'\preceq_dD'_0\prec_dD'$. Also, by (2), some sufficiently high tail
of $R$ is strictly smaller in the debt order than $D'$.
The level $w$ of $\U$ is positive, possibly infinite.

First suppose $\alpha_m<\infty$ along a subsequence. By (3), for large
$m$ a non-null tail $W=R\cap\{t>\alpha_{m'}\}$, $m'$ sufficiently large,
lies in $\cred_i(C_m)$ and makes $(A',W)$ strictly AC. Since $T_y$
misses $R$, whereas $A'\prec_dD'$ and $(D',T_y)$ is SI, (IB) gives
\[
 V+\debt_i(A')+\cred_i(T_y)\prec_i V
 \prec_i V+\debt_i(A')+\cred_i(W),\qquad W\prec_cT_y.
\]
If $w=\infty$, the tail $W$ belongs to $\U$, a contradiction.
If $w<\infty$, take the initial $m$ large enough that $\beta_m>w+1$,
then choose $m'$ and $W$ as above, and choose
$B\in\U$, $B\subseteq R\cap\debt_i(C_m)$, with
$a(B)+b(B)<b(W)$. Then $W\prec_cT_y\prec_c B$.
Replace the debt on $B$ by credit and silence $W$.
At $i$ this is a strict improvement: credit on $B$ is better than credit
on $W$, even before the old debt on $B$ is counted. The change of
$\Lambda_j$ is $a(B)+b(B)-b(W)<0$. Contradiction.

Otherwise $C_m=E_j(\beta_m)$ along a subsequence. Fix such a model $E$
with $\beta_m>w+1$ if $w<\infty$. Choose a non-null
$Z\subseteq R\cap\text{silence}(E)$ in a bounded interval of a high tail, so that
$Z\prec_dD'$. If $w<\infty$, choose
$B\in\U$, $B\subseteq R\cap\debt_i(E)$, with $a(B)+b(B)<a(Z)$.
Changing debt to credit on $B$ and adding debt on $Z$ lowers
$\Lambda_j$. Competitiveness and (IB), followed by (G) to remove the
initial debt from the left side, give
$V\preceq_iV+\debt_i(Z)+\cred_i(B)$.
If $w=\infty$, instead choose $B\in\U$ in a tail beyond $Z$, with
$b(B)<a(Z)$; adding credit on $B$ and debt on $Z$ gives the same
comparison. In either case, $T_y\prec_c B$ and $Z\prec_dD'$ imply
\[
 V\prec_i V+\debt_i(Z)+\cred_i(T_y)
 \prec_i V+\debt_i(D')+\cred_i(T_y)\prec_i V,
\]
the final contradiction. Every common context used here is silent on
the sets being compared: $T_y$ is disjoint from $R$ and from $D'$.
\end{proof}

\begin{theorem}\label{thm:allfine}
For every finite $n\ge3$, under complete (G)+(IB) rankings and (S),
if one world is fine, all worlds are fine. Consequently every coherent
model whose credal set contains three vertices is dominated by a selected
coherent model. In particular the vacuous model is dominated.
\end{theorem}
\begin{proof}
Let $k$ be fine and fix $i\ne k$. Edge transfer gives
$D_{ik}\in\J^i_A\cap\K^i_T$ for all non-null $A,T$.
It also makes every world $j\ne k$ fine on $D_{jk}$.

If $i$'s debt order were not fine, choose an ultrafilter $\U$ avoiding
some $\J^i_A$. It cannot contain $D_{ik}$. Since the finitely many
$D_{ij}$ cover $H'$, it contains some $D_{ij}$ with $j\notin\{i,k\}$,
hence $R=D_{ij}\setminus D_{ik}$. Lemma~\ref{lem:remaining}(a)
excludes this. Thus $i$'s debt order is fine.

For credit the same argument supplies $\U$ containing such an $R$.
Input (6), on $[e_i,e_k]$ with fine partner $k$, gives the SI pair
$(D',T_y)$ with $T_y\preceq^i_cT$ and $T_y\subseteq D_{ik}$, so
$T_y$ misses $R$. Lemma~\ref{lem:remaining}(b) excludes $\U$.
Thus $i$ is fine. No undominatedness assumption on $V$ has been used.

FSE's local-extension proposition on $H'$ now gives a strictly positive
finitely additive weak extension at every world. The CBOI obstruction,
in its one-sided form recorded in CCF, applies for every $n\ge3$ and
dominates each coherent model containing three vertices. The KKM
consequence of (S) makes a selected coherent dominator available.
\end{proof}

\section{A fine edge, and a world with an additive extension}

Dominating $V$ at three worlds need not dominate it at the other worlds.
For the next two arguments we therefore use narrower regions than the
triple-sign regions above. Define
\[
 O_v=\{g:g_vg_j<0\text{ for every }j\ne v\}.
\]
Only $v$ disagrees with all the other worlds on $O_v$. Write $a_v$ for
the advice there that is bad at $v$ and good at every other world.
The regions $O_v$ are disjoint and non-null. They need not cover $H'$.
For distinct $u,v,w$, all nonempty open conditions on the two ratios
used below have non-null intersections with the appropriate $O_v$.
Any nonempty modification of $V$ using only $a_u,a_v,a_w$ on these
three regions strictly improves $V$ at every world outside
$\{u,v,w\}$. This observation makes the three-world dominance
arguments apply without discarding the remaining worlds.

\begin{theorem}\label{thm:fineedge}
For every finite $n\ge3$, under complete (G)+(IB) rankings and (S),
one fine edge forces $V$ to be dominated.
\end{theorem}
\begin{proof}
Suppose the edge $[e_v,e_w]$ is fine and $V$ is undominated.
Fix $u\notin\{v,w\}$. We check FSE's fine-edge argument with
$O_u,O_v,O_w$ in place of its three sign regions.
For a two-sided selection $E$ on $[e_u,e_v]$, put
\[
 T_0=O_v\cap\{t^u_v>\alpha_E\},\qquad
 B_0=O_u\cap\{t^u_v<\beta_E\}.
\]
Both are non-null. On them $E$ gives respectively $a_v$ and $a_u$.

First, for every non-null $B\subseteq B_0$ there is a non-null high
tail $Z_B\subseteq O_w$ such that
$V+a_u(B)+a_w(Z)\prec_wV$ for every non-null $Z\subseteq Z_B$.
This is FSE's Lemma B1: obtain an SI pair inside $D_{vw}$ from its
local extension; replace its credit by $B$ using a low-cone yardstick,
and its debt by a sufficiently high tail inside $O_w$ using a debt
yardstick. The tail in $O_w$ is non-null, and $B\subseteq O_u$ misses
$D_{vw}$, just as required in that proof.

Second, suppose $V+a_v(T)+a_u(B)\preceq_uV$ for some $T\subseteq T_0$.
Partition $T$ into finitely many pieces, each of whose debt at $v$ is
strictly outweighed by credit on $Z_B$, using $v$'s fine local
extension on $D_{vw}$. Add the pieces in order. There is a first step
at which the displayed comparison holds at $u$. At the preceding
step it holds at $v$: competitiveness of $E$, with its advice silenced
on the changed sets, and (IB) give this alternative.
Adding $a_w(Z_B)$ at the first step now strictly improves $V$ at $u$,
at $v$, and at $w$, exactly as in FSE's Lemma B2. It strictly improves
at every remaining world too, since each of $a_u,a_v,a_w$ is good
there. This contradicts undominatedness. Consequently
\[
 V\prec_uV+a_v(T)+a_u(B)
 \quad(T\subseteq T_0,\; B\subseteq B_0\text{ non-null}).
\]

Third, $T_0$ is credit-null at $u$. Otherwise an SI pair with credit
inside $T_0$ has its debt replaceable by a non-null set
$B_0\cap\{t^u_w>b\}$ for sufficiently large $b$, by the debt yardstick.
This would contradict the last display (FSE's Lemma B3). But
Theorem~\ref{thm:LH}, proved independently below, excludes a non-null
credit-null region. Hence no two-sided selection exists on
$[e_u,e_v]$: that edge has the endpoint property and is fine, by
FSE's endpoint-to-fineness lemma.

This holds for every $u\ne v,w$. Together with the original fine
edge, all edges at $v$ are fine. Input (4) and the finite cover
$H'=\bigcup_{u\ne v}D_{vu}$ make $v$ a fine world.
Theorem~\ref{thm:allfine} contradicts the assumed undominatedness of $V$.
\end{proof}

\begin{theorem}\label{thm:oneextension}
For every finite $n\ge3$, under complete (G)+(IB) rankings and (S),
if one world has a strictly positive finitely additive extension of
CBOI's form for comparisons from coherent undominated models to
regular models, then $V$ is dominated.
\end{theorem}
\begin{proof}
Suppose $V$ is undominated and world $v$ has the stated extension,
with charges $\nu_1,\nu_2$ on $H'$. Fix distinct $u,w\ne v$ and
suppose a selection $C$ on $[e_u,e_v]$ is two-sided. Define $T_0,B_0$
as above. Let $Z_m\downarrow$ be common pure sets for the charges,
with $\mu(Z_m)\to0$; outside $Z_m$ they have positive densities.

The preliminary edge results in FSE's tame-world argument already
hold for every finite $n\ge3$. They say that a one-sided extension at
$v$ supplies the partner bound on $D_{vw}$, makes $w$ fine on
$D_{vw}\setminus Z_m$ relative to every non-null yardstick, and, for
large $m$, supplies an SI pair there with both parts non-null.
The almost agreeing debt and credit probabilities on this region
admit arbitrarily fine finite partitions.

Choose $m_1$ so that $T_0\setminus Z_{m_1}$ and
$B_0\setminus Z_{m_1}$ are non-null. By the density of $\nu_2$ there,
choose a non-null $T\subseteq T_0\setminus Z_{m_1}$ with
\[
 \nu_2(T)<\nu_1(B_0\setminus Z_{m_1}).
\]
FSE's third-world exchange (Step 2 of its chain-extension theorem)
then gives a non-null $S\subseteq O_w$ such that
\[
 V+a_v(T)+a_w(S)\prec_w V.                          \tag{*}
\]
Here are the placement details in the present dimension. Take
$m\ge m_1$ with $O_w\setminus Z_m$ non-null, and an SI pair
$(A_0,T'')$ inside $D_{vw}\setminus Z_m$. If needed shrink $A_0$ and
$T$ so that they are disjoint, preserving their non-nullity and the
last inequality. Divide the SI pair into equally many pieces so that
every credit piece is smaller in its almost agreeing probability than
$T$. At least one piece is SI, since otherwise composing the AC
pieces would make the original pair AC. Replace that piece's credit
by $T$, and its debt by a sufficiently small non-null part
$S\subseteq O_w\setminus Z_m$ in the debt probability. This proves
(*). All replacements have disjoint debt and credit parts. These are
exactly the operations in FSE's Step 2; no partition of $H'$ into
three regions is used.

Fill $B_0\setminus Z_{m_1}$ continuously in measure by an increasing
family $\Sigma_s$, $0\le s\le1$. The function
$\lambda(s)=\nu_2(T)-\nu_1(\Sigma_s)$ is continuous, positive at $0$
and negative at $1$. Choose $s$ with $0<\lambda(s)<\nu_1(S)$, and set
\[
 X'=V+a_w(S),\qquad X=X'+a_v(T)+a_u(\Sigma_s),
\]
and let $Y$ be $C$ silenced on $T\cup\Sigma_s$.
At world $v$ the loss differences are
\[
 L_v(X)-L_v(V)=\lambda(s)-\nu_1(S)<0,\qquad
 L_v(Y)-L_v(C)=-\lambda(s)<0.
\]
The one-sided extension, applied at the coherent undominated $V,C$,
gives $V\not\preceq_vX$ and $C\not\preceq_vY$.
Also $X\prec_wV$ by (*) and (G), and $X\prec_jV$ at every world
$j\notin\{u,v,w\}$, because its advice is good there.

Undominatedness and completeness therefore force $V\preceq_uX$.
Competitiveness of $C$ on $\{u,v\}$ forces $C\preceq_uY$; (IB)
turns this into $X\preceq_uX'$. But $X'\prec_uV$ by (G), a
contradiction. Thus every edge $[e_u,e_v]$ has the endpoint property
and is fine. Their finite union makes $v$ fine, and
Theorem~\ref{thm:allfine} gives the final contradiction.
\end{proof}

\section{Neither kind of advice is negligible}

A region $R$ is \emph{credit-null} at $v$ if every disjoint pair $(A,T)$
with $A$ non-null and $T\subseteq R$ is AC. It is
\emph{debt-infinitesimal} if every disjoint pair $(A,T)$ with
$A\subseteq R$ non-null and $T$ non-null is SI. Let $\mathrm{CN}_v$
and $\mathrm{DI}_v$ be these hereditary families, as in FSE.

\begin{theorem}\label{thm:LH}
For every finite $n\ge3$, under complete (G)+(IB) rankings and (S),
no world has a
non-null credit-null region or a non-null debt-infinitesimal region.
No assumption about $V$ being undominated is needed.
\end{theorem}
\begin{proof}
Here are the dimension checks for FSE's low-cone-heaviness proof.
For any three distinct $v,w,u$, replace its three regions by
\[
 P_v^{vwu}=D_{vw}\cap D_{vu},\quad
 P_w^{vwu}=D_{wv}\cap D_{wu},\quad
 P_u^{vwu}=D_{uv}\cap D_{uw}.
\]
They are disjoint regions where the named world disagrees with the
other two. They need not cover $H'$; no step below requires that.
Every prescribed nonempty open condition on the two ratios in one of
these regions has positive measure.

First, a world cannot have both kinds of non-null region: disjoint
non-null parts give a pair required to be both SI and AC (FSE's exclusion
lemma). Second, FSE's cones-and-tails lemma holds on \emph{each}
specified edge. Its proof uses only the head and tail bounds in input
(2), splitting the original region in two to make the context disjoint.
Thus a credit-null region produces a credit-null low cone on each edge,
and a debt-infinitesimal region produces a debt-infinitesimal high tail
on each edge. When two edges are needed, take the smaller of the two
head levels, or the larger of the two tail levels.

Third, the endpoint-property argument excluding credit-null low cones
uses a single edge and is unchanged. Fourth, FSE's transfer from
credit-nullity at $v$ to debt-infinitesimality at $w$ is valid with
$P_w^{vwu}$ in place of $P_w$. Its auxiliary set is
\[
 Q\subseteq P_w^{vwu}\cap\{t^v_w<\beta_E\}
                   \cap\{t^w_u<a'\},
\]
which is non-null, in the debt of $E$ at $v$, and no larger in $w$'s
credit order than the desired credit yardstick, by (2). Silencing
$Q$ and the credit-null part in $E$, and applying competitiveness on
$\{v,w\}$, gives precisely the same SI comparison at $w$.

Finally, FSE's transfer from $\mathrm{DI}_w\ne0$ to
$\mathrm{CN}_v\ne0$ also holds. Its first case uses only input (3).
In its second case, a selection near $D_v$ has thresholds
$(\alpha,0)$ at $v$. Take the candidate credit-null region
$T'=P_w^{vwu}\cap\{t^v_w<\alpha\}$.
For any non-null debt yardstick $A$, (2) supplies a high $t^v_u$-tail
no larger than $A$; its intersection with
\[
 P_v^{vwu}\cap\{\alpha/4<t^v_w<\alpha/2\}
\]
is a non-null set $S\preceq^v_d A$, silent in the selection and
disjoint from $T'$. Adding debt at $w$ on a part of $T'$ and credit at
$w$ on $S$ strictly improves the selection at $w$. Competitiveness
and (IB) give the AC comparison at $v$, and replacement of $S$ by $A$
finishes. This verifies the only place where FSE's yardstick must be
placed in a specified triple-sign region.

Now suppose $\mathrm{CN}_v\ne0$ and choose distinct $w,u\ne v$.
Both edges at $v$ fail the endpoint property; take two-sided selections
$E$ on $[e_v,e_w]$ and $E'$ on $[e_v,e_u]$. A high $t^v_w$-tail
intersected with a credit-null low $t^v_u$-cone inside $P_v^{vwu}$
is non-null and lies in the credit of $E$; the analogous intersection
lies in the credit of $E'$. The transfers just checked give
$\mathrm{DI}_w\ne0$ and $\mathrm{DI}_u\ne0$, and then
$\mathrm{CN}_u\ne0$, contradicting exclusion at $u$.
Thus every $\mathrm{CN}_v$ is null. A non-null $\mathrm{DI}_w$ would
produce a non-null $\mathrm{CN}_v$ at any other world, so those too
are excluded.
\end{proof}

\begin{corollary}
In Theorem~\ref{thm:allfine}, it suffices that one world satisfy any of
FSE's sufficient conditions for fineness: germ continuity, either stated
form of conic continuity, or the endpoint property on all its incident
edges. The other worlds need no additional continuity assumption.
\end{corollary}

\section{What this does not settle}

None of these theorems forces a world or an edge to be fine. The general complete-ranking
problem is still open, already for three worlds. The first theorem
settles FSE's higher-dimensional propagation question;
Theorems~\ref{thm:fineedge}--\ref{thm:oneextension} extend its fine-edge
and one-sided-extension obstructions, and Theorem~\ref{thm:LH}
settles its corresponding higher-dimensional low-cone question.
They do not extend the two-additive-world theorem without completeness
to arbitrary $n$.
\end{document}
